\documentclass[../main.tex]{subfiles}
\begin{document}
\setlist[itemize]{topsep=2pt,itemsep=1pt,parsep=0pt,partopsep=0pt}
\setlist[itemize,2]{topsep=1pt,itemsep=0pt,parsep=0pt,partopsep=0pt}
\setlength{\parskip}{3pt}

\chapter{Week 4 — Validation and Sustainability}
\label{ch:week4}

\textit{Snapshot and part summaries: see the overview on page~\pageref{ov:week4}.}

\section{Part 1 — Verification and Validation}
\label{sec:vv}

\begin{itemize}
  \item \keyterm{Verification}: the method was applied correctly and is suitable for the topic — internal, process-facing check.
  \item \keyterm{Validation}: the outcomes are true and supported by strong evidence — external, reality-facing check.
  \item Both are required for credibility, professional usability, and academic reproducibility.
  \item Industry: verification = quality control against specifications (internal); validation = quality assurance of fitness for purpose with users (external).
  \item Thesis guidance: ask in advance how methods will be verified and how the goal can be proven true; design the validation exercise into the schedule rather than assuming there will be time for it.
\end{itemize}

\section{Part 2 — Engineering for One Planet}
\label{sec:eop}

\begin{itemize}
  \item \keyterm{EOP framework}: planning must consider environmental, social, cultural, and economic implications across three axes.
  \item \keyterm{Value of research}: forecast near- and long-term cost and value; ask whether there is potential harm, whether the work can be built upon, and whether it contributes to science or industry.
  \item \keyterm{Reproducibility}: signals transparency and rigour; barriers are complexity, technology change, human error, and IP rights — good data management mitigates the latter two (\autoref{sec:dmp}).
  \item \keyterm{Resource management} — digital: favour local over server, local storage over cloud, download frequently accessed documents, print only what will be read three or more times.
  \item Resource management — methodology: verify that the target accuracy is realistic, whether a less intensive alternative exists, and whether the sample size is justified against the objective (\autoref{sec:sample-types}).
\end{itemize}

\section{Reading Material}
\label{sec:reading4}

\begin{partbox}
The Week~4 reading, \textit{Conducting Sustainable Research}, expands EOP into a working method, linking the value axis to literature-gap analysis and the reproducibility axis to concrete data-management habits.
\end{partbox}

\begin{itemize}
  \item Tie research value to a literature gap judged on its environmental, social, cultural, and economic implications.
  \item Use scripted analysis software over graphical tools to eliminate human error and ensure identical repetition.
  \item Use descriptive, self-explanatory file and variable names with version dates, in non-proprietary formats (text, csv).
  \item Keep an uncorrected original data file and document the full workflow from raw data to results.
  \item Troubleshoot locally, compress transferred files, and delete server files once work is complete.
  \item Garbage in equals garbage out: justify sample size against the objective.
\end{itemize}

\end{document}
